%% kapa: exam=1 round=21 year=2010 period=1 form=A subject=civil_law subject_ko=민법 seq=33 number=33
%% source: Q-Net 감정평가사 기출문제 (cst00309) artlSeq=5024181 file=extracted/1008/5024181/2010년 제21회 감정평가사 1차시험 기출문제/1교시.hwp
%% conversion: kordoc 4.24.51 -> tools/gen1.py (auto); review=proofread-sonnet; flags=none
\begin{question}{33}
甲은 2004년 4월 2일 10시에 가출하였고, 甲의 생사를 알 수 없다. 다음 설명으로 옳지 않은 것은? (다툼이 있으면 판례에 의함)
\begin{choices}
  \item 甲에게 선순위 상속인이 있는 경우, 후순위 상속인은 실종선고를 청구할 수 없다.
  \item 2010년 2월 13일 甲에 대한 실종선고가 확정된 경우, 甲은 2009년 4월 2일 24시에 사망한 것으로 본다.
  \item 甲의 부동산을 상속한 乙이 그 부동산을 丙에게 양도한 후 甲이 생환한 경우, 생환만으로 그 부동산의 소유권이 甲에게 당연히 복귀하는 것은 아니다.
  \item 甲에 대한 실종선고가 취소되기 위해서는 이해관계인 등의 청구가 있어야 하지만, 공시최고절차를 거쳐야 하는 것은 아니다.
  \item 甲의 생환으로 인한 실종선고취소는 甲의 상속인 乙이 실종기간 만료 후 선고 전에 선의로 한 법률행위의 효력에는 영향을 미치지 않는다.
\end{choices}
\end{question}
\answer{5}
